% Springer LNCS conference paper draft (HWID 2026 style)
\documentclass[runningheads]{llncs}
\usepackage[T1]{fontenc}
\usepackage{graphicx}
\usepackage{url}

\begin{document}

\title{Treating Workers Like Guided Missile Pigeons: Cybernetics, Behaviorism, and Symbolic Systems of Control}
\titlerunning{Treating Workers Like Guided Missile Pigeons}

\author{Watson Hartsoe}
\authorrunning{W. Hartsoe}
\institute{Digital Media, Georgia Institute of Technology\\
\email{whartsoe3@gatech.edu}}

\maketitle

\begin{abstract}
Industry 5.0 is commonly framed as a human-centric paradigm that moves beyond productivity to worker wellbeing and augmentation. This paper argues the opposite: the Fifth Industrial Revolution (5IR) does not transcend earlier paradigms of behavioral control, but internalizes them. Tracing a genealogy from Norbert Wiener's cybernetics and B.F. Skinner's Project Pigeon (later Project ORCON) to contemporary workplace analytics, we argue that modern ``human-in-the-loop'' systems reuse a mid-century architecture in which living organisms are integrated into cybernetic control loops as servomechanisms. We introduce the concept of the \emph{Thick Machine}: a system whose operational efficacy is inseparable from cultural legitimacy. Project ORCON failed not simply because of technical limits, but because pigeon-guided warfare was semantically intolerable to military command. Contemporary systems such as Reinforcement Learning from Human Feedback (RLHF) and affective workplace surveillance succeed by solving that symbolic problem: they rebrand conditioning as care, and optimization as empathy. We conclude by outlining six conditions for preserving \emph{symbolic sovereignty} in human-machine teams, including the right to illegibility, semantic autonomy, and refusal of enforced interpretation.

\keywords{Cybernetics \and Behaviorism \and Symbolic Anthropology \and Industry 5.0 \and Digital Twins \and ORCON \and Thick Description}
\end{abstract}

\section{Introduction: The Semiotic Refusal}

Project Pigeon was not only a technical experiment. It was a cultural scandal.

In 1943, B.F. Skinner proposed a guidance system for the ``Pelican'' glide bomb. At the time, reliable electronic guidance remained limited, vulnerable, or expensive. Skinner's proposal was biological: train pigeons to peck at the image of a target on a projection screen; translate pecks into steering corrections through a servo mechanism; use operant conditioning to keep the birds locked on target \cite{skinner1960pigeons,capshew1993engineering}.

The system worked well enough in controlled tests to establish a principle. A living organism could be placed inside a cybernetic loop as an error-correcting component. The bird's internal intent was irrelevant. The loop only required stable sensory discrimination, conditioned response, and mechanical transduction.

Yet the project was canceled. Official explanations emphasized practicality and military confidence in competing electronic systems, but Skinner's own account notes the committee's restrained amusement and skepticism toward the image of pigeons guiding munitions \cite{skinner1960pigeons}. The failure was not merely engineering. It was semiotic. The architecture functioned; the meaning of the architecture was intolerable.

This paper argues that contemporary ``human-centric'' AI systems solve precisely that symbolic problem. The worker is now integrated into feedback systems as an organic control component, but the arrangement is rendered culturally acceptable through a layer of interpretation: wellness, coaching, resilience, collaboration, and augmentation. What was once visibly conditioning is now framed as care.

We use \emph{Project ORCON} (Organic Control) as the historical reference point for this architecture. We then extend the analysis through Clifford Geertz's concept of thick description \cite{geertz1973interpretation} to show how contemporary systems extract not only labor but interpretive authority over the worker's interior life. The result is a form of control that is behaviorist in mechanism and symbolic in presentation.

\section{Six Observations}

The argument proceeds through six observations.

\begin{enumerate}
\item \textbf{The cybernetic-behaviorist synthesis treats organisms as servomechanisms.} Internal states are secondary to the control problem. What matters is reliable input-output behavior under feedback.

\item \textbf{RLHF operationalizes culture as a Skinner box.} Human value judgments are translated into reward signals that shape model policy, reducing moral and cultural norms to statistical reinforcement.

\item \textbf{The locus of extraction shifts from labor power to semiotic surplus.} The system increasingly captures the worker's capacity to define the meaning of their own stress, attention, emotion, and flourishing.

\item \textbf{The Thick Machine succeeds where the Thin Machine failed.} Project ORCON exposed conditioning too directly. Contemporary systems wrap control in the language of empathy, safety, and optimization.

\item \textbf{The operator-in-the-loop is often a rhetorical illusion.} The human is positioned as a supervisor, but functionally serves as an error-correcting component for edge cases and model failures.

\item \textbf{Symbolic sovereignty is the right to illegibility and interpretation.} Human-centric design requires the right to withhold interiority from algorithmic oversight and contest machine interpretation.
\end{enumerate}

\section{Cybernetics and Behaviorism as a Control Grammar}

Norbert Wiener's wartime work on anti-aircraft prediction helped establish a foundational cybernetic ontology: the pilot, aircraft, and weapon could be modeled as a coupled control system governed by feedback, prediction, and correction \cite{wiener1948cybernetics}. Under combat stress, the pilot becomes legible to the predictor not as a sovereign subject but as a patterning organism.

Cybernetics therefore provided a language for treating human and machine behavior as functionally comparable at the level of control: both regulate deviation from target states through loops of sensing, acting, and correction.

Skinner's behaviorism supplied the complementary programming logic. Radical behaviorism bracketed interior mental states and focused on the shaping of observable behavior via environmental contingencies and reinforcement schedules \cite{skinner1953science}. The Skinner box demonstrated that complex behavioral repertoires could be engineered through carefully calibrated reward and punishment.

When combined, cybernetics and behaviorism produce a totalizing grammar of control. The organism becomes a black-box component whose outputs can be optimized if the environment, reward schedule, and feedback loop are properly designed. Project Pigeon is not an eccentric exception to this synthesis. It is one of its clearest prototypes.

\section{Project ORCON and the Thick Machine}

Skinner's pigeons were selected for practical reasons: visual acuity, trainability, and rapid response \cite{skinner1960pigeons}. Their pecks were transduced into steering corrections through a servo linkage. The architecture was elegant. The bird's eye served as sensor; the beak and neck movements served as actuators; grain served as reinforcement.

Project ORCON reveals an important asymmetry between technical function and cultural legitimacy. In a narrow engineering sense, the system asks a simple question: can an organism reliably generate corrective behavior when exposed to a moving image target? In a military-symbolic sense, the system asks a different question: what does it mean for command, responsibility, and sovereignty if a farm bird occupies the control loop of a modern weapon?

This distinction motivates the concept of the \emph{Thick Machine}. A Thick Machine is not merely a technical apparatus. It is a sociotechnical system whose usability depends on the cultural meaning of its components, roles, and visible rituals. Project ORCON failed as a Thick Machine because the arrangement violated the military's image of rational, masculine, electromechanical command \cite{schultzfigueroa2019project,wlodarczyk2020beyond}. The pigeon-guided bomb was operationally plausible and semantically obscene.

Geertz's distinction between a twitch and a wink clarifies the issue \cite{geertz1973interpretation}. To the control engineer, a peck is a kinematic event in a feedback loop. To the institution, the same event carries symbolic meaning. The peck is read not just as motion, but as a statement about agency, dignity, and technological legitimacy.

The cancellation of ORCON, then, can be read as a case of semiotic refusal: the institution rejected not simply a weapon, but an intolerable meaning.

\section{From Thin Control to Thick Control}

Early cybernetic systems largely operated in what we can call a thin mode of description. They optimized trajectories, timings, and error tolerances while ignoring the symbolic density of human interpretation. A target was a coordinate. A correction was a signal. A response was a measurable output.

Contemporary AI systems cannot remain thin. They increasingly operate in semantic environments: customer interactions, workplace communication, moderation, coaching, education, recommendation, and assessment. In these settings, the system must not only regulate behavior but interpret signs of intention, affect, and social meaning.

This creates a historical inversion. Mid-century behaviorism dismissed thick human interpretation as methodologically inconvenient. Platform-era AI depends on the extraction and operationalization of thick human judgments at scale. The system must continuously translate meaning into machine-readable reward and control variables.

The result is not the abandonment of behaviorism, but its expansion. The mechanism remains thin. The interface becomes thick.

\section{RLHF as a Cultural Skinner Box}

Reinforcement Learning from Human Feedback (RLHF) is the most visible example of this expansion. In the RLHF pipeline, models generate outputs; human evaluators rank or score them; these rankings are compressed into reward signals; the model is optimized to maximize expected reward \cite{sutton2018reinforcement}. The loop is mathematically continuous with operant conditioning.

From a control perspective, RLHF solves an alignment problem by building a learned approximation of acceptable human response. From a cultural perspective, it performs a conversion: it turns ethics, style, civility, and normativity into trainable gradients.

This is why RLHF is best understood as a cultural Skinner box. The model does not acquire ethics in a philosophical sense. It acquires a policy that navigates a statistical terrain of approval. It learns the topography of reward, not the ontology of the good.

The same logic applies recursively to human labor in annotation and quality control pipelines. Human raters are themselves scored, calibrated, and normalized to produce consistent labels. They are conditioned to condition the model. Human cultural judgment is transformed into a managed component of a larger optimization system.

\section{Industry 5.0 and the Sanitized ORCON}

Industry 5.0 is publicly framed as human-centric, resilient, and sustainable \cite{breque2021industry}. This framing is politically effective because it positions advanced automation as a corrective to the dehumanization associated with earlier industrial paradigms. The worker returns to the center.

But a closer reading of human-in-the-loop systems suggests a different interpretation. The worker does not simply return as an autonomous craft subject. The worker returns as a highly instrumented organism embedded in adaptive feedback architectures that monitor, predict, and modulate behavior.

The modern ORCON solves the symbolic problem that doomed Project Pigeon. It does not place the worker in a visible nose cone. It places them inside an ``intelligent workspace.'' It does not offer grain. It offers engagement scores, coaching, resilience analytics, and wellness interventions. The conditioning remains; the semiotics change.

\subsection{Operator Digital Twins and Affective Surveillance}

The Operator Digital Twin (ODT) extends digital-twin logic from machinery to the worker's body and behavior. Wearables, computer vision, productivity telemetry, and linguistic analytics can be combined to produce a dynamic model of fatigue, attention, posture, affect, or stress risk. In contact center settings, commercial products already market emotion AI as real-time coaching and wellbeing support \cite{cxtoday2022cogito}.

Here the ORCON crosses from thin behavioral measurement into thick interpretive capture. Physiological signals (heart rate variability, speech cadence, micro-pauses) are treated as proxies for internal states such as engagement, resilience, empathy, or burnout risk. The system maps biometric twitch to symbolic wink.

This mapping is never neutral. It installs a machine interpretation of the worker's interiority and links that interpretation to intervention: slower pace, different tasks, alerts, nudges, coaching prompts, breaks, or escalations. The worker is told the system is caring for them; functionally, the system is maintaining a biological component inside a production loop.

In this sense, the ODT is not a straightforward augmentation tool. It is a coercive prosthetic relation: the worker and model are coupled, but the terms of coupling are set by institutional objectives and opaque interpretive models rather than by the worker's own account of their condition.

\section{The Capture of Flourishing}

The normative language of ``flourishing,'' ``wellbeing,'' and ``human-centricity'' becomes technically unstable when absorbed into cybernetic systems of optimization. Flourishing is a thick moral concept shaped by history, philosophy, and social struggle. It is not reducible to a single measurable state.

Yet operational systems require measurable variables. Once flourishing enters the ORCON, it is translated into bounded indicators: acceptable heart-rate variability ranges, stable attention signatures, non-volatile sentiment, sustainable task throughput, reduced safety incidents. The system then optimizes toward these proxies.

This is a capture operation. The institution gains the power to define the measurable signs of a good life at work, and then to enforce those signs through adaptive environment design. If the system needs calm, calm becomes wellness. If it needs intensified output under narrow thresholds of distress, resilience becomes wellness. The worker's own interpretation of exhaustion or refusal is displaced by dashboard truth.

The locus of extraction therefore shifts. Beyond surplus labor, the system extracts \emph{semiotic surplus}: the worker's capacity to interpret and narrate the meaning of their own embodied state.

\section{The Semiotic Paradox}

Modern ORCON systems rely on thick human meaning but process it through thin control mechanisms. They require richly contextual human judgments to train reward models, and they require culturally situated interpretations to map behavior to affect. Yet at the point of operation, these meanings are compressed into scalar rewards, risk scores, or threshold triggers.

The paradox is structural. The system depends on interpretive richness to function, but it can only govern through reduction. It therefore oscillates between semantic dependence and semantic violence.

This also explains why contemporary systems are often experienced as empathic while acting coercively. The appearance of understanding comes from high-resolution pattern capture and adaptive response. But true empathy would require acknowledging the subject's irreducible alterity and their right to contest the system's reading. ORCON-style systems instead substitute homeostatic intervention for interpretation. The worker is not understood; the worker is regulated.

\section{Counter-Cybernetics and Symbolic Sovereignty}

If the modern ORCON governs by stabilizing machine interpretations of human behavior, resistance cannot be reduced to privacy in the narrow sense. The core issue is symbolic sovereignty: the right to retain interpretive standing over one's own embodied and emotional life.

We therefore propose six conditions for symbolic sovereignty in human-machine work systems.

\begin{enumerate}
\item \textbf{The Right to Illegibility.} Workers must be able to selectively withhold biometric, behavioral, or affective data from the system without punitive outcomes.

\item \textbf{The Right to Semantic Autonomy.} Workers must be able to contest and override algorithmic interpretations of their behavior, emotion, and risk state.

\item \textbf{The Right to Interiority.} Physiological and emotional data should not be captured as a free input to optimization systems; use must be consensual, bounded, and compensable.

\item \textbf{The Right to Agency.} Human-in-the-loop systems must preserve genuine human direction rather than reduce workers to exception handlers for machine control regimes.

\item \textbf{The Right to Friction and Refusal.} Workers must be able to reject ``seamless'' interventions and wellness nudges without being classified as irrational, unsafe, or noncompliant.

\item \textbf{The Right to Interpretation.} Final authority over the meaning of fatigue, stress, morale, and flourishing must remain with human persons and communities, not with dashboards or reward models.
\end{enumerate}

These conditions do not reject automation. They establish a boundary: the system may assist, but it may not monopolize meaning.

\section{Conclusion}

Project Pigeon was dismissed as eccentric, yet it anticipated a central problem of contemporary AI systems: how to integrate living beings into control loops while preserving institutional legitimacy.

Skinner proved a narrow technical point. If one controls environment, stimuli, and reinforcement, one can produce reliable corrective action from an organism without requiring its understanding of the larger system. Industry 5.0 extends that principle into the workplace but solves the earlier semiotic crisis through a new interface layer. The loop is no longer presented as conditioning. It is presented as support.

This is the triumph of the Thick Machine. Behavioral control persists, but now depends on the capture of symbolic meaning. The worker is not only measured but interpreted; not only optimized but narrated. The infrastructure of care becomes the infrastructure of control.

The critical design question for human-machine systems is therefore not whether humans remain ``in the loop.'' It is who defines the loop's meanings, who controls the interpretive layer, and whether workers retain the right to become partially unreadable to the systems that manage them.

The pigeons did not disappear. They were professionalized.

% ---- Bibliography ----
\begin{thebibliography}{99}

\bibitem{breque2021industry}
Breque, M., De Nul, L., Petridis, A.: Industry 5.0: Towards a sustainable, human-centric and resilient European industry. European Commission, Directorate-General for Research and Innovation, Brussels (2021)

\bibitem{capshew1993engineering}
Capshew, J.H.: Engineering behavior: Project Pigeon, World War II, and the conditioning of B.F. Skinner. Technology and Culture \textbf{34}(4), 835--857 (1993)

\bibitem{cxtoday2022cogito}
Isobel, M.: Cogito updates its emotion AI software for real-time coaching. CX Today (2022), \url{https://www.cxtoday.com/contact-centre/cogito-updates-its-emotion-ai-software/}

\bibitem{geertz1973interpretation}
Geertz, C.: Thick Description: Toward an Interpretive Theory of Culture. In: The Interpretation of Cultures: Selected Essays, pp. 3--30. Basic Books, New York (1973)

\bibitem{kellogg2020algorithms}
Kellogg, K.C., Valentine, M.A., Christin, A.: Algorithms at work: The new contested terrain of control. Academy of Management Annals \textbf{14}(1), 366--410 (2020)

\bibitem{schultzfigueroa2019project}
Schultz-Figueroa, B.: Project Pigeon: Rendering the war animal through optical technology. JCMS: Journal of Cinema and Media Studies \textbf{58}(4), 92--111 (2019)

\bibitem{skinner1953science}
Skinner, B.F.: Science and Human Behavior. Macmillan, New York (1953)

\bibitem{skinner1960pigeons}
Skinner, B.F.: Pigeons in a pelican. American Psychologist \textbf{15}(1), 28--37 (1960)

\bibitem{sutton2018reinforcement}
Sutton, R.S., Barto, A.G.: Reinforcement Learning: An Introduction, 2nd edn. MIT Press, Cambridge (2018)

\bibitem{wiener1948cybernetics}
Wiener, N.: Cybernetics: Or Control and Communication in the Animal and the Machine. MIT Press, Cambridge (1948)

\bibitem{wlodarczyk2020beyond}
Wlodarczyk, J.: Beyond bizarre: Nature, culture and the spectacular failure of B.F. Skinner's pigeon-guided missiles. Polish Journal for American Studies \textbf{14}, 7--20 (2020)

\end{thebibliography}

\end{document}
